
Group Relative Policy Optimization (GRPO) has attracted interest as a training method for code-generating language models because it replaces the learned critic of PPO with a group-relative normalization scheme, reducing implementation complexity while retaining execution feedback as the reward signal. However, the group-relative formulation carries a structural precondition: the advantage function

$$A_i = \frac{r_i - \text{mean}(r_{\text{group}})}{\text{std}(r_{\text{group}})}$$

is defined only when std(r_group) > 0. When all G completions in a group receive identical reward—most commonly, when all fail and receive reward zero—the advantage is undefined and the gradient contribution is zero. This degenerate case is not an edge case for 7B-class models on competitive programming tasks: models of this scale solve competitive programming problems at near-zero rates, so the majority of training groups are degenerate throughout training.

The failure mode is not visible in standard training diagnostics. Loss curves may appear stable. Periodic benchmark evaluations may show marginal or noisy changes. The collapse becomes apparent only when per-step advantage variance is measured directly.

In this paper, we address the following question: does task granularity—specifically, whether execution feedback is binary (function-level) or allows partial credit (repository-level)—determine whether GRPO advantage variance is in a viable or degenerate regime for 7B-class code models?

We conduct a controlled experiment comparing two conditions: function-level GRPO training on APPS and CodeContests (binary execution reward), and repository-level GRPO training on SWE-bench Verified (partial credit reward based on file-path overlap). All other factors are held constant: the same model (CodeLlama-7b-Instruct-hf), the same GRPO configuration (G=8, B=4), and the same training duration (120 steps per condition). We measure per-step advantage variance at every training step.

We additionally implement and smoke-test a training infrastructure for a 4-condition curriculum GRPO experiment on APPS+CodeContests, as a prerequisite for a full behavioral test of whether easy-to-hard curriculum ordering improves pass@1 on HumanEval+. The full behavioral experiment was not executed.

The remainder of this paper is organized as follows. Section 2 reviews related work. Section 3 describes the methodology. Section 4 describes the experimental setup. Section 5 reports results. Section 6 discusses implications and limitations. Section 7 concludes.
